\documentclass[../DoAn.tex]{subfiles}
\begin{document}

\section{Thiết kế kiến trúc}
\label{section:4.1}

\subsection{Lựa chọn kiến trúc phần mềm}
\label{subsection:4.1.1}

Hệ thống ALIBARBIE được tổ chức theo kiến trúc \textbf{ba tầng} (3-tier
architecture) kết hợp mô hình \textbf{MVC} (Model--View--Controller) phía máy
chủ. Kiến trúc ba tầng phân chia hệ thống thành ba lớp độc lập:

\begin{itemize}
  \item \textbf{Tầng trình bày (Presentation Tier):} Giao diện React.js chạy
        trên trình duyệt, giao tiếp với máy chủ qua REST API.
  \item \textbf{Tầng nghiệp vụ (Business Logic Tier):} Máy chủ Node.js +
        Express.js xử lý toàn bộ logic nghiệp vụ và cung cấp API.
  \item \textbf{Tầng dữ liệu (Data Tier):} MongoDB lưu trữ dữ liệu, truy cập
        thông qua Mongoose ODM.
\end{itemize}

Trong tầng nghiệp vụ, mô hình MVC được áp dụng cụ thể như sau:

\begin{itemize}
  \item \textbf{Model (M):} 17 schema Mongoose định nghĩa cấu trúc dữ liệu và
        ánh xạ trực tiếp với các collection MongoDB (\texttt{ModelUser},
        \texttt{ModelProducts}, \break \texttt{ModelOrder}, v.v.).
  \item \textbf{Controller (C):} 15 lớp controller xử lý nghiệp vụ, nhận yêu
        cầu từ tầng Route, tương tác với Model và trả về JSON response
        (\texttt{ControllerUser},\break \texttt{ControllerAdmin},
        \texttt{ControllerReturn}, v.v.).
  \item \textbf{View (V):} Do hệ thống xây dựng theo mô hình API thuần túy
        (headless), tầng View không nằm ở phía máy chủ mà được thay thế bởi
        các component React phía client---mỗi component nhận dữ liệu JSON từ
        API và tự render giao diện.
\end{itemize}

Giữa Route và Controller, hệ thống bổ sung \textbf{tầng Middleware} để xử lý
xác thực JWT (\texttt{verifyToken}) và kiểm tra quyền theo vai trò
(\texttt{verifyRole}), tạo ra luồng xử lý bốn bước:
\[
  \text{Route} \;\longrightarrow\; \text{Middleware (JWT + RBAC)} \;\longrightarrow\;
  \text{Controller} \;\longrightarrow\; \text{Model}
\]

Kiến trúc này đáp ứng yêu cầu bảo trì và mở rộng từng module độc lập như đã
xác định ở mục \ref{section:2.4}: mỗi nhóm chức năng (đơn hàng, hoàn hàng,
vận chuyển, phân quyền, v.v.) có controller và model riêng, không ảnh hưởng
lẫn nhau khi thay đổi.

\subsection{Thiết kế tổng quan}
\label{subsection:4.1.2}

Hình \ref{fig:package-diagram} minh họa biểu đồ phụ thuộc gói của hệ thống.
Hệ thống được chia thành hai siêu gói: \textbf{Client} (phía trình duyệt) và
\textbf{Server} (phía máy chủ), giao tiếp với nhau qua REST API và Socket.io.

\begin{figure}[H]
    \centering
    \includegraphics[width=0.9\textwidth]{Hinhve/package_diagram.png}
    \caption{Biểu đồ phụ thuộc gói của hệ thống ALIBARBIE}
    \label{fig:package-diagram}
\end{figure}

\textbf{Phía Server}, các gói được tổ chức theo tầng từ trên xuống:

\begin{itemize}
  \item \texttt{route} --- Tiếp nhận HTTP request, phân phối đến middleware và
        controller phù hợp. Gói này phụ thuộc vào \texttt{controller} và
        \texttt{middleware}.
  \item \texttt{middleware} --- Chứa \texttt{ControllerJWT} thực hiện xác thực
        token và kiểm tra quyền trước khi chuyển tiếp xuống controller. Phụ
        thuộc vào \texttt{model} để tra cứu vai trò tùy chỉnh.
  \item \texttt{controller} --- Xử lý toàn bộ nghiệp vụ; phụ thuộc vào
        \texttt{model} và \texttt{utils}.
  \item \texttt{model} --- Định nghĩa schema Mongoose; phụ thuộc vào
        \texttt{config} (kết nối database).
  \item \texttt{config} --- Cấu hình kết nối MongoDB; không phụ thuộc gói nào
        khác.
  \item \texttt{utils} --- Các tiện ích dùng chung (ChatBot, auditLog); phụ
        thuộc vào \texttt{model}.
\end{itemize}

\textbf{Phía Client}, các gói được tổ chức theo tầng từ trên xuống:

\begin{itemize}
  \item \texttt{Pages} --- Các trang ứng dụng (storefront và trang quản trị);
        phụ thuộc vào \texttt{Layouts}, \texttt{components}, \texttt{redux} và
        \texttt{config}.
  \item \texttt{Layouts} --- Header, Footer, Loading---các thành phần khung
        trang; phụ thuộc vào \texttt{components}.
  \item \texttt{components} --- Các component tái sử dụng nhỏ (CascadingAddressSelect,
        VietnamCitySelect); phụ thuộc vào \texttt{config}.
  \item \texttt{redux} --- Store, reducer, actions---quản lý trạng thái giỏ
        hàng toàn cục; không phụ thuộc gói giao diện nào.
  \item \texttt{contexts} --- PermissionContext cung cấp thông tin quyền hạn
        cho \break component con; phụ thuộc vào \texttt{config}.
  \item \texttt{Route} --- Định nghĩa bảng định tuyến publicRoutes và
        privateRoute; phụ thuộc vào \texttt{Pages}.
  \item \texttt{config} --- URL gốc của API; không phụ thuộc gói nào khác.
\end{itemize}

Các gói tầng dưới không phụ thuộc ngược lên tầng trên (ví dụ \texttt{model}
không phụ thuộc \texttt{controller}), đảm bảo nguyên tắc thiết kế phân tầng
một chiều.

\subsection{Thiết kế chi tiết gói}
\label{subsection:4.1.3}

Hình \ref{fig:pkg-server-controller} mô tả quan hệ giữa các lớp trong nhóm
gói \texttt{route}--\texttt{middleware}--\texttt{controller} phía server. Đây
là nhóm trung tâm hiện thực hóa toàn bộ luồng xử lý API.

\begin{figure}[H]
    \centering
    \includegraphics[width=0.9\textwidth]{Hinhve/pkg_server.png}
    \caption{Thiết kế gói server: route -- middleware -- controller}
    \label{fig:pkg-server-controller}
\end{figure}

Lớp \texttt{index} (route) \textit{phụ thuộc} vào sáu lớp route con
(\texttt{UserRoutes}, \texttt{ProductsRoutes}, \texttt{AdminRoutes},
\texttt{PaymentsRoutes}, \texttt{WebRoutes}) và gọi trực tiếp một số
controller dùng cho endpoint công khai. \texttt{AdminRoutes} có quan hệ
\textit{phụ thuộc} vào \texttt{ControllerJWT} (middleware kiểm tra JWT và
RBAC) trước khi chuyển tiếp đến các controller tương ứng. Mỗi controller
(\texttt{ControllerUser}, \texttt{ControllerAdmin}, v.v.) có quan hệ
\textit{kết hợp} với một hoặc nhiều lớp Model để thực hiện truy vấn dữ liệu.

Hình \ref{fig:pkg-client} mô tả quan hệ giữa các lớp trong gói
\texttt{Pages/Admin} phía client---gói phức tạp nhất do quản lý nhiều tính
năng quản trị.

\begin{figure}[H]
    \centering
    \includegraphics[width=0.9\textwidth]{Hinhve/pkg_client.png}
    \caption{Thiết kế gói client: Pages/Admin}
    \label{fig:pkg-client}
\end{figure}

\texttt{DefaultLayout} (component cha) có quan hệ \textit{hợp thành} với \break
\texttt{HeaderAdmin}, \texttt{SlideBar} và \texttt{HomePageAdmin}.
\texttt{HomePageAdmin} có quan hệ \textit{hợp thành} với các sub-page:
\texttt{Dashboard}, \texttt{Products}, \break \texttt{OrderProducts},
\texttt{Customers}, \texttt{Categories}, \texttt{Coupons}, \texttt{Returns},
\texttt{Shipping}, \texttt{RolePermission}, \texttt{History}, v.v.
\texttt{PermissionContext} có quan hệ \textit{phụ thuộc} đến
\texttt{DefaultLayout} để cung cấp thông tin quyền hạn cho toàn bộ cây
component bên trong.

%%%%%%%%%%%%%%%%%%%%%%%%%%%%%%%%%%%%%%%%%%
\section{Thiết kế chi tiết}
\label{section:4.2}

\subsection{Thiết kế giao diện}
\label{subsection:4.2.1}

\subsubsection*{Đặc tả màn hình mục tiêu}

Hệ thống hướng đến hai nhóm người dùng với thiết bị khác nhau:

\begin{itemize}
  \item \textbf{Khách hàng mua sắm:} Truy cập chủ yếu qua điện thoại di động
        (màn hình 360px--414px) và máy tính để bàn (1366px--1920px). Giao
        diện phải tương thích đầy đủ trên cả hai nhóm.
  \item \textbf{Quản trị viên/nhân viên:} Truy cập qua máy tính để bàn,
        độ phân giải tối thiểu 1280~×~720 px.
\end{itemize}

Hệ thống hỗ trợ đầy đủ màu 24-bit (16 triệu màu), font chữ UTF-8 để hiển
thị tiếng Việt không dấu lỗi.

\subsubsection*{Chuẩn hóa giao diện}

\begin{itemize}
  \item \textbf{Hệ thống lưới:} Sử dụng grid 12 cột của Bootstrap 5 với các
        breakpoint chuẩn (sm: 576px, md: 768px, lg: 992px, xl: 1200px).
  \item \textbf{Bảng màu:} Màu chủ đạo hồng pastel (\texttt{\#ff6b9d}) cho
        thương hiệu \break ALIBARBIE; màu nền trắng (\texttt{\#ffffff}); màu chữ
        chính xám đậm (\texttt{\#212529}); màu cảnh báo/lỗi đỏ
        (\texttt{\#dc3545}); màu thành công xanh (\texttt{\#198754}).
  \item \textbf{Nút bấm (Button):} Nút chính (primary) nền hồng, chữ trắng;
        nút phụ \break (secondary) viền xám; nút nguy hiểm (danger) nền đỏ dùng cho
        thao tác xóa. Kích thước tối thiểu vùng chạm 44~×~44 px trên thiết bị
        di động.
  \item \textbf{Thông điệp phản hồi:} Sử dụng \texttt{react-toastify}---thông
        báo xuất hiện góc trên phải màn hình, tự động ẩn sau 3 giây. Màu
        xanh = thành công; màu đỏ = lỗi; màu vàng = cảnh báo.
  \item \textbf{Form nhập liệu:} Nhãn (label) hiển thị phía trên input; thông
        báo lỗi validation hiển thị ngay dưới trường tương ứng bằng chữ đỏ
        nhỏ.
  \item \textbf{Ảnh sản phẩm:} Lazy loading thông qua \texttt{react-lazyload}
        để giảm thời gian tải trang; tỷ lệ khung hình cố định 1:1 (vuông)
        cho thumbnail.
  \item \textbf{Phân trang:} Sử dụng \texttt{react-paginate} với số mục mỗi
        trang có thể cấu hình; số trang hiển thị tối đa 5 nút.
\end{itemize}

\subsubsection*{Minh họa thiết kế các giao diện chính}

Hình \ref{fig:ui-home} minh họa thiết kế trang chủ storefront với bố cục
Header (logo, thanh tìm kiếm, giỏ hàng) -- Banner slideshow -- Lưới sản phẩm
nổi bật -- Footer.

\begin{figure}[H]
    \centering
    \includegraphics[width=0.9\textwidth]{Hinhve/ui_home.png}
    \caption{Thiết kế giao diện trang chủ}
    \label{fig:ui-home}
\end{figure}

Hình \ref{fig:ui-checkout} minh họa thiết kế màn hình thanh toán với bố cục
hai cột: cột trái hiển thị form địa chỉ giao hàng và phương thức thanh toán;
cột phải tóm tắt giỏ hàng và tổng tiền (bao gồm phí vận chuyển và giảm giá
coupon).

\begin{figure}[H]
    \centering
    \includegraphics[width=0.9\textwidth]{Hinhve/ui_checkout.png}
    \caption{Thiết kế giao diện trang thanh toán}
    \label{fig:ui-checkout}
\end{figure}

Hình \ref{fig:ui-admin} minh họa thiết kế trang quản trị với bố cục sidebar
bên trái (danh sách menu theo quyền hạn) và vùng nội dung bên phải thay đổi
theo menu được chọn.

\begin{figure}[H]
    \centering
    \includegraphics[width=0.9\textwidth]{Hinhve/ui_admin.png}
    \caption{Thiết kế giao diện trang quản trị}
    \label{fig:ui-admin}
\end{figure}

\subsection{Thiết kế lớp}
\label{subsection:4.2.2}

\subsubsection*{Các lớp chủ đạo}

Bảng \ref{tab:class-user} mô tả thiết kế lớp \texttt{ControllerUser}---lớp
quản lý toàn bộ vòng đời tài khoản người dùng từ đăng ký, đăng nhập đến cập
nhật hồ sơ.

\begin{table}[H]
\centering
\caption{Thiết kế lớp \texttt{ControllerUser}}
\label{tab:class-user}
\scriptsize
\begin{tabular}{|p{4cm}|p{9cm}|}
\hline
\multicolumn{2}{|c|}{\textbf{ControllerUser}} \\
\hline
\multicolumn{2}{|l|}{\textit{Thuộc tính}} \\
\hline
(không có thuộc tính) & Lớp sử dụng các Model và service bên ngoài \\
\hline
\multicolumn{2}{|l|}{\textit{Phương thức}} \\
\hline
\texttt{Register(req, res)} & Kiểm tra email tồn tại, băm mật khẩu bằng bcrypt (saltRounds=10), lưu \texttt{ModelUser}, ghi audit log \\
\hline
\texttt{Login(req, res)} & Tìm user theo email, so sánh mật khẩu bcrypt, tạo JWT (maxAge 1h hoặc 30 ngày), set cookie \\
\hline
\texttt{GetUser(req, res)} & Giải mã cookie JWT, trả về thông tin user hiện tại \\
\hline
\texttt{ChangePass(req, res)} & Xác thực mật khẩu cũ, băm mật khẩu mới, cập nhật MongoDB \\
\hline
\texttt{EditProfile(req, res)} & Cập nhật fullname, phone, avatar của user \\
\hline
\texttt{Logout(req, res)} & Xóa cookie Token \\
\hline
\texttt{PostComment(req, res)} & Lưu bình luận sản phẩm vào \texttt{ModelComments} \\
\hline
\texttt{GetDataOrder(req, res)} & Lấy danh sách đơn hàng theo email từ \texttt{ModelOrder} \\
\hline
\end{tabular}
\end{table}

Bảng \ref{tab:class-return} mô tả thiết kế lớp \texttt{ControllerReturn}---lớp
quản lý quy trình yêu cầu hoàn hàng nhiều bước.

\begin{table}[H]
\centering
\caption{Thiết kế lớp \texttt{ControllerReturn}}
\label{tab:class-return}
\scriptsize
\begin{tabular}{|p{4cm}|p{9cm}|}
\hline
\multicolumn{2}{|c|}{\textbf{ControllerReturn}} \\
\hline
\multicolumn{2}{|l|}{\textit{Phương thức}} \\
\hline
\texttt{CreateRequest(req, res)} & Tạo yêu cầu hoàn hàng mới: kiểm tra đơn hàng tồn tại, lưu \break \texttt{ModelReturnRequest}, ghi audit log \\
\hline
\texttt{GetMyRequests(req, res)} & Lấy danh sách yêu cầu của khách hàng theo email \\
\hline
\texttt{GetAllRequests(req, res)} & Lấy toàn bộ yêu cầu (dành cho nhân viên/quản lý) \\
\hline
\texttt{GetRequestById(req, res)} & Lấy chi tiết một yêu cầu theo ID \\
\hline
\texttt{UpdateStatus(req, res)} & Cập nhật trạng thái yêu cầu (PENDING~→~CONTACTING~→~...) \\
\hline
\texttt{ApproveReturn(req, res)} & Duyệt hoàn hàng: cập nhật status = APPROVED, gửi email thông báo \\
\hline
\texttt{RejectReturn(req, res)} & Từ chối: cập nhật status = REJECTED, lưu lý do, gửi email \\
\hline
\texttt{ProcessRefund(req, res)} & Xử lý hoàn tiền: cộng số dư vào tài khoản, status = REFUNDED \\
\hline
\end{tabular}
\end{table}

Bảng \ref{tab:class-admin} mô tả lớp \texttt{ControllerAdmin}---lớp trung
tâm cho giao diện quản trị.

\begin{table}[H]
\centering
\caption{Thiết kế lớp \texttt{ControllerAdmin} (các phương thức chính)}
\label{tab:class-admin}
\scriptsize
\begin{tabular}{|p{4.5cm}|p{8.5cm}|}
\hline
\multicolumn{2}{|c|}{\textbf{ControllerAdmin}} \\
\hline
\multicolumn{2}{|l|}{\textit{Phương thức}} \\
\hline
\texttt{GetDataOrder(req, res)} & Lấy toàn bộ đơn hàng, populate thông tin sản phẩm từ \break \texttt{ModelOrderItem} \\
\hline
\texttt{EditOrder(req, res)} & Cập nhật trạng thái đơn hàng; nếu hoàn tất gửi email xác nhận \\
\hline
\texttt{AddProduct(req, res)} & Thêm sản phẩm mới vào \texttt{ModelProducts}, hỗ trợ upload ảnh \\
\hline
\texttt{EditProduct(req, res)} & Cập nhật thông tin sản phẩm, ghi audit log \\
\hline
\texttt{DeleteProduct(req, res)} & Xóa mềm hoặc xóa sản phẩm, ghi audit log \\
\hline
\texttt{GetUser(req, res)} & Trả về danh sách toàn bộ tài khoản \\
\hline
\texttt{UpdateUserRole(req, res)} & Thay đổi vai trò nhân viên, ghi audit log \\
\hline
\texttt{GetDataAuth(req, res)} & Trả về thông tin tài khoản đang đăng nhập cùng danh sách quyền \\
\hline
\end{tabular}
\end{table}

\subsubsection*{Biểu đồ trình tự}

\textbf{Use case 1: Đăng nhập hệ thống.}
Hình \ref{fig:seq-login} mô tả luồng truyền thông điệp khi người dùng thực
hiện đăng nhập.

\begin{figure}[H]
    \centering
    \includegraphics[width=0.9\textwidth]{Hinhve/seq_login.png}
    \caption{Biểu đồ trình tự: Đăng nhập hệ thống}
    \label{fig:seq-login}
\end{figure}

\begin{figure}[H]
    \centering
    \includegraphics[width=0.9\textwidth]{Hinhve/seq_register.png}
    \caption{Biểu đồ trình tự: Đăng ký người dùng}
    \label{fig:seq-register}
\end{figure}

Luồng chính: (1) Client gửi \texttt{POST /api/login} với \{email, password,
rememberMe\}. (2) \texttt{ControllerUser.Login} gọi
 \break \texttt{ModelUser.findOne(\{email\})} để tìm tài khoản. (3) Nếu tìm thấy,
gọi \texttt{bcrypt.compare()} để kiểm tra mật khẩu. (4) Nếu khớp, tạo JWT
bằng \texttt{jwt.sign()} và set cookie \texttt{Token}. (5) Trả về JSON với
\{message, role, fullname, avatar\}. (6) Client lưu thông tin vào
\texttt{localStorage} và chuyển hướng trang.

\textbf{Use case 2: Đặt hàng và thanh toán VNPay.}
Hình \ref{fig:seq-order} mô tả luồng truyền thông điệp khi khách hàng hoàn
tất đặt hàng qua cổng VNPay.

\begin{figure}[H]
    \centering
    \includegraphics[width=0.9\textwidth]{Hinhve/seq_dat_hang.png}
    \caption{Biểu đồ trình tự: Đặt hàng}
    \label{fig:seq-order}
\end{figure}

\begin{figure}[H]
    \centering
    \includegraphics[width=0.9\textwidth]{Hinhve/seq_thanh_toan.png}
    \caption{Biểu đồ trình tự: Thanh toán}
    \label{fig:seq-payment}
\end{figure}

Luồng chính: (1) Client gửi \texttt{POST /api/payment} với thông tin đơn
hàng và địa chỉ giao hàng. (2) \texttt{PaymentController} tính phí vận
chuyển qua \break \texttt{ControllerShipping}, kiểm tra coupon (nếu có), tạo bản ghi
 \break \texttt{ModelOrder} và \texttt{ModelOrderItem}. (3) Server tạo URL thanh toán
\break VNPay có chữ ký HMAC-SHA512 và trả về URL này. (4) Client chuyển hướng
trình duyệt đến cổng VNPay. (5) Sau khi khách hàng thanh toán, VNPay gọi
 \break callback \texttt{GET /vnpay-return}. (6) Server xác thực chữ ký và cập nhật
\break \texttt{statusPayment = true} trong \texttt{ModelOrder}. (7) Server gửi email
xác nhận qua \break \texttt{Nodemailer}, xóa giỏ hàng và trả về kết quả.

\textbf{Use case 3: Xử lý yêu cầu hoàn hàng.}
Các hình dưới đây mô tả luồng từ khi khách hàng tạo yêu cầu đến khi
quản lý hoàn tiền.

\begin{figure}[H]
    \centering
    % SV bổ sung sequence diagram hoàn hàng tại đây
    \includegraphics[width=0.9\textwidth]{Hinhve/seq_hoan_hang_1.png}
    \caption{Biểu đồ trình tự: Xử lý yêu cầu hoàn hàng: Kiểm tra điều kiện và tạo yêu cầu trả hàng
}
    \label{fig:seq-return-1-4}
\end{figure}

\begin{figure}[H]
    \centering
    % SV bổ sung sequence diagram hoàn hàng tại đây
    \includegraphics[width=0.9\textwidth]{Hinhve/seq_hoan_hang_2.png}
    \caption{Biểu đồ trình tự: Xử lý yêu cầu hoàn hàng: Liên hệ khách hàng và khách gửi hàng
}
    \label{fig:seq-return-2-4}
\end{figure}

\begin{figure}[H]
    \centering
    % SV bổ sung sequence diagram hoàn hàng tại đây
    \includegraphics[width=0.9\textwidth]{Hinhve/seq_hoan_hang_3.png}
    \caption{Biểu đồ trình tự: Xử lý yêu cầu hoàn hàng: Duyệt yêu cầu
}
    \label{fig:seq-return-3-4}
\end{figure}

\begin{figure}[H]
    \centering
    % SV bổ sung sequence diagram hoàn hàng tại đây
    \includegraphics[width=0.9\textwidth]{Hinhve/seq_hoan_hang_4.png}
    \caption{Biểu đồ trình tự: Xử lý yêu cầu hoàn hàng: Khách hàng theo dõi trạng thái trả hàng
}
    \label{fig:seq-return-4-4}
\end{figure}

\begin{figure}[H]
    \centering
    % SV bổ sung sequence diagram hoàn hàng tại đây
    \includegraphics[width=0.9\textwidth]{Hinhve/seq_hoan_tien.png}
    \caption{Biểu đồ trình tự: Xử lý hoàn tiền
}
    \label{fig:seq-refund}
\end{figure}

Luồng chính: (1) Khách hàng gửi \break \texttt{POST /api/return-request} với \{order\_id,
reason, description, images\}. (2) \texttt{ControllerReturn.CreateRequest} kiểm
tra đơn hàng tồn tại, lưu \break \texttt{ModelReturnRequest} với status
\texttt{PENDING\_REVIEW}. (3) Nhân viên cập nhật trạng thái lần lượt
(\texttt{CONTACTING} → \texttt{WAITING\_ITEM} → \texttt{ITEM\_RECEIVED}).
(4) Quản lý gọi \texttt{POST /api/admin/return-requests/:id/approve}.
(5) \texttt{ControllerReturn.ApproveReturn} cập nhật status = \texttt{APPROVED}
và gửi email thông báo. (6) Quản lý tiếp tục gọi endpoint \texttt{refund};
server cộng \texttt{refund\_amount} vào trường \texttt{surplus} của
\texttt{ModelUser} và đặt \break status = \texttt{REFUNDED}.

\subsection{Thiết kế cơ sở dữ liệu}
\label{subsection:4.2.3}

Do sử dụng MongoDB (NoSQL), thiết kế cơ sở dữ liệu được mô tả qua sơ đồ
thực thể liên kết (E-R diagram) điều chỉnh cho mô hình document, kết hợp với
mô tả cấu trúc từng collection.

\subsubsection*{Sơ đồ E-R}

Hình \ref{fig:er-diagram} trình bày sơ đồ E-R của hệ thống với 17 collection.

\begin{figure}[H]
    \centering
    \includegraphics[width=\textwidth]{Hinhve/er_diagram.png}
    \caption{Sơ đồ E-R hệ thống ALIBARBIE}
    \label{fig:er-diagram}
\end{figure}

Các quan hệ chính giữa các collection:

\begin{itemize}
  \item \texttt{products} (1) --- (N) \texttt{product\_variants}: Mỗi sản phẩm
        có nhiều biến thể màu sắc, mỗi biến thể chứa danh sách kích cỡ nhúng
        (embedded \texttt{sizes}).
  \item \texttt{order} (1) --- (N) \texttt{orderItem}: Mỗi đơn hàng chứa
        nhiều dòng sản phẩm; \texttt{orderItem} lưu tham chiếu \texttt{order\_id}
        và \texttt{product\_id}.
  \item \texttt{order} (1) --- (0..1) \texttt{return\_request}: Một đơn hàng
        có thể có tối đa một yêu cầu hoàn hàng.
  \item \texttt{user} (1) --- (N) \texttt{order}: Mỗi người dùng có thể có
        nhiều đơn hàng; liên kết qua trường \texttt{email}.
  \item \texttt{products} (N) --- (1) \texttt{category}: Mỗi sản phẩm thuộc
        một danh mục; liên kết qua \texttt{category\_id} (ObjectId reference).
  \item \texttt{role} (1) --- (N) \texttt{user}: Mỗi vai trò có thể gán cho
        nhiều người dùng; liên kết qua trường \texttt{role} (String).
  \item \texttt{audit\_log} ghi nhận hành động của \texttt{user} lên các đối
        tượng bất kỳ qua \texttt{target\_type} và \texttt{target\_id} (mô hình
        polymorphic reference).
\end{itemize}

\subsubsection*{Thiết kế các collection chính}

Bảng \ref{tab:col-user} đến Bảng \ref{tab:col-return} mô tả cấu trúc các
collection quan trọng.

\begin{table}[H]
\centering
\caption{Collection \texttt{user}}
\label{tab:col-user}
\scriptsize
\begin{tabular}{|p{3.5cm}|p{3cm}|p{6cm}|}
\hline
\textbf{Trường} & \textbf{Kiểu dữ liệu} & \textbf{Mô tả} \\
\hline
\texttt{\_id} & ObjectId & Khóa chính (tự sinh) \\
\hline
\texttt{fullname} & String & Họ và tên \\
\hline
\texttt{email} & String & Email (duy nhất, dùng làm định danh) \\
\hline
\texttt{password} & String & Mật khẩu đã băm bcrypt \\
\hline
\texttt{role} & String (enum) & Vai trò: admin / manager / staff / user \\
\hline
\texttt{isAdmin} & Boolean & Cờ quản trị tối cao \\
\hline
\texttt{phone} & Number & Số điện thoại \\
\hline
\texttt{surplus} & Number & Số dư tài khoản (hoàn tiền) \\
\hline
\texttt{avatar} & String & Đường dẫn ảnh đại diện \\
\hline
\texttt{created\_at} & Date & Thời điểm tạo \\
\hline
\end{tabular}
\end{table}

\begin{table}[H]
\centering
\caption{Collection \texttt{products}}
\label{tab:col-products}
\scriptsize
\begin{tabular}{|p{3.5cm}|p{3cm}|p{6cm}|}
\hline
\textbf{Trường} & \textbf{Kiểu dữ liệu} & \textbf{Mô tả} \\
\hline
\texttt{\_id} & ObjectId & Khóa chính \\
\hline
\texttt{nameProducts} & String & Tên sản phẩm \\
\hline
\texttt{priceNew} & Number & Giá bán hiện tại \\
\hline
\texttt{priceOld} & Number & Giá gốc (để hiển thị khuyến mãi) \\
\hline
\texttt{category\_id} & ObjectId (ref) & Tham chiếu danh mục \\
\hline
\texttt{img} & String & Ảnh đại diện sản phẩm \\
\hline
\texttt{images} & [String] & Danh sách ảnh chi tiết \\
\hline
\texttt{stock\_quantity} & Number & Tồn kho tổng \\
\hline
\texttt{free\_shipping} & Boolean & Miễn phí vận chuyển \\
\hline
\texttt{return\_days} & Number & Số ngày được phép hoàn trả \\
\hline
\texttt{rating\_avg} & Number & Điểm đánh giá trung bình \\
\hline
\end{tabular}
\end{table}

\begin{table}[H]
\centering
\caption{Collection \texttt{product\_variants} (biến thể sản phẩm)}
\label{tab:col-variant}
\scriptsize
\begin{tabular}{|p{3.5cm}|p{3cm}|p{6cm}|}
\hline
\textbf{Trường} & \textbf{Kiểu dữ liệu} & \textbf{Mô tả} \\
\hline
\texttt{\_id} & ObjectId & Khóa chính \\
\hline
\texttt{product\_id} & ObjectId (ref) & Tham chiếu sản phẩm cha \\
\hline
\texttt{color} & String & Tên màu \\
\hline
\texttt{color\_hex} & String & Mã màu hex (hiển thị ô màu) \\
\hline
\texttt{img} & String & Ảnh biến thể màu sắc \\
\hline
\texttt{sizes} & [SizeSchema] & Danh sách kích cỡ nhúng (size, stock, price\_adjustment) \\
\hline
\end{tabular}
\end{table}

\begin{table}[H]
\centering
\caption{Collection \texttt{order} và \texttt{orderItem}}
\label{tab:col-order}
\scriptsize
\begin{tabular}{|p{3.5cm}|p{3cm}|p{6cm}|}
\hline
\textbf{Trường} & \textbf{Kiểu dữ liệu} & \textbf{Mô tả} \\
\hline
\multicolumn{3}{|c|}{\textit{order}} \\
\hline
\texttt{email} & String & Email khách hàng đặt hàng \\
\hline
\texttt{sumPrice} & Number & Tổng tiền đơn hàng (đã gồm ship, giảm giá) \\
\hline
\texttt{statusOrder} & Boolean & Trạng thái giao hàng \\
\hline
\texttt{statusPayment} & Boolean & Trạng thái thanh toán \\
\hline
\texttt{payment\_method} & String (enum) & Phương thức: vnpay / cod \\
\hline
\texttt{vnp\_txn\_ref} & String & Mã giao dịch VNPay \\
\hline
\texttt{has\_return\_request} & Boolean & Cờ đánh dấu đã có yêu cầu hoàn hàng \\
\hline
\multicolumn{3}{|c|}{\textit{orderItem}} \\
\hline
\texttt{order\_id} & ObjectId (ref) & Tham chiếu đơn hàng \\
\hline
\texttt{product\_id} & ObjectId (ref) & Tham chiếu sản phẩm \\
\hline
\texttt{nameProduct} & String & Tên sản phẩm tại thời điểm đặt \\
\hline
\texttt{quantity} & Number & Số lượng \\
\hline
\texttt{price} & Number & Đơn giá tại thời điểm đặt \\
\hline
\end{tabular}
\end{table}

\begin{table}[H]
\centering
\caption{Collection \texttt{return\_request}}
\label{tab:col-return}
\scriptsize
\begin{tabular}{|p{3.5cm}|p{3cm}|p{6cm}|}
\hline
\textbf{Trường} & \textbf{Kiểu dữ liệu} & \textbf{Mô tả} \\
\hline
\texttt{order\_id} & ObjectId (ref) & Tham chiếu đơn hàng \\
\hline
\texttt{customer\_email} & String & Email khách hàng yêu cầu \\
\hline
\texttt{reason} & String & Lý do hoàn hàng \\
\hline
\texttt{images} & [String] & Ảnh minh chứng \\
\hline
\texttt{status} & String (enum) & PENDING\_REVIEW / CONTACTING \break /  WAITING\_ITEM / ITEM\_RECEIVED \break / APPROVED / REJECTED / REFUNDED \\
\hline
\texttt{refund\_amount} & Number & Số tiền hoàn trả \\
\hline
\texttt{reject\_reason} & String & Lý do từ chối (nếu có) \\
\hline
\texttt{modified\_by} & String & Email nhân viên xử lý cuối \\
\hline
\end{tabular}
\end{table}

Collection \texttt{audit\_log} được lập chỉ mục tổng hợp trên
\texttt{(actor\_email, created\_at)}, \texttt{(target\_type, target\_id,
created\_at)} và \break \texttt{action\_code} để đảm bảo hiệu năng truy vấn nhật ký
theo nhiều chiều phân tích (mục \ref{subsection:2.3.4}).

%%%%%%%%%%%%%%%%%%%%%%%%%%%%%%%%%%%%%%%%%%
\section{Xây dựng ứng dụng}
\label{section:4.3}

\subsection{Thư viện và công cụ sử dụng}
\label{subsection:4.3.1}

Bảng \ref{tab:tools-server} và Bảng \ref{tab:tools-client} liệt kê các công
cụ và thư viện sử dụng để phát triển hệ thống.

\begin{table}[H]
\centering
\caption{Công cụ và thư viện phía Server}
\label{tab:tools-server}
\scriptsize
\begin{tabular}{|p{4.5cm}|p{3.5cm}|p{5cm}|}
\hline
\textbf{Mục đích} & \textbf{Công cụ / Thư viện} & \textbf{Phiên bản} \\
\hline
IDE lập trình & Visual Studio Code & 1.89+ \\
\hline
Nền tảng thực thi & Node.js & 20.x LTS \\
\hline
Framework web & Express.js & 4.19.2 \\
\hline
Hệ quản trị CSDL & MongoDB & 7.0 \\
\hline
ODM (Object-Document Mapper) & Mongoose & 8.3.2 \\
\hline
Xác thực JWT & jsonwebtoken & 9.0.2 \\
\hline
Băm mật khẩu & bcrypt & 5.1.1 \\
\hline
Giao tiếp thời gian thực & Socket.io & 4.7.5 \\
\hline
Tích hợp thanh toán & vnpay & 1.5.0 \\
\hline
Gửi email & Nodemailer & 6.9.13 \\
\hline
OAuth2 Google & googleapis & 135.0.0 \\
\hline
Upload file & Multer & 1.4.5-lts.1 \\
\hline
HTTP client & Axios & 1.7.2 \\
\hline
Dev server tự động reload & Nodemon & 3.1.0 \\
\hline
Kiểm thử API & Postman & v11+ \\
\hline
Quản lý biến môi trường & dotenv & 16.4.5 \\
\hline
\end{tabular}
\end{table}

\begin{table}[H]
\centering
\caption{Công cụ và thư viện phía Client}
\label{tab:tools-client}
\scriptsize
\begin{tabular}{|p{4.5cm}|p{3.5cm}|p{5cm}|}
\hline
\textbf{Mục đích} & \textbf{Công cụ / Thư viện} & \textbf{Phiên bản} \\
\hline
Thư viện giao diện & React & 18.2.0 \\
\hline
Quản lý trạng thái & Redux Toolkit & 2.2.3 \\
\hline
Kết nối Redux-React & React-Redux & 9.1.1 \\
\hline
Định tuyến client & React Router DOM & 6.22.3 \\
\hline
CSS Framework & Bootstrap & 5.3.3 \\
\hline
Component Bootstrap & React-Bootstrap & 2.10.2 \\
\hline
Biểu đồ doanh thu & Chart.js / react-chartjs-2 & 4.4.2 / 5.2.0 \\
\hline
Biểu đồ doanh thu (phụ) & Recharts & 3.8.1 \\
\hline
Icon & FontAwesome & 6.x \\
\hline
Thông báo toast & react-toastify & 10.0.5 \\
\hline
Phân trang & react-paginate & 8.2.0 \\
\hline
Lazy load ảnh & react-lazyload & 3.2.1 \\
\hline
Slider banner & react-slick & 0.30.2 \\
\hline
Socket client & socket.io-client & 4.7.5 \\
\hline
HTTP client & Axios & 1.6.8 \\
\hline
Xử lý CSS & Sass & 1.75.0 \\
\hline
\end{tabular}
\end{table}

\subsection{Kết quả đạt được}
\label{subsection:4.3.2}

Sản phẩm đạt được gồm hai ứng dụng độc lập chạy song song:

\begin{itemize}
  \item \textbf{Máy chủ API} (\texttt{server/}): Ứng dụng Node.js cung cấp
        hơn 60 API endpoint REST và dịch vụ Socket.io, phục vụ cả giao diện
        mua sắm lẫn trang quản trị.
  \item \textbf{Ứng dụng web} (\texttt{client/}): SPA React với hai giao diện
        tách biệt: storefront cho khách hàng (tuyến \texttt{/}) và trang quản
        trị cho nội bộ (tuyến \texttt{/admin}).
\end{itemize}

Bảng \ref{tab:stats} thống kê thông tin về mã nguồn.

\begin{table}[H]
\centering
\caption{Thống kê mã nguồn hệ thống ALIBARBIE}
\label{tab:stats}
\scriptsize
\begin{tabular}{|p{5cm}|p{3cm}|p{3cm}|p{2cm}|}
\hline
\textbf{Thành phần} & \textbf{Số file} & \textbf{Số dòng code} & \textbf{Ngôn ngữ} \\
\hline
Server (Node.js/Express) & 46 & 3.535 & JavaScript \\
\hline
Client (React) & 64 & 8.494 & JavaScript/JSX \\
\hline
\textbf{Tổng cộng} & \textbf{110} & \textbf{12.029} & \\
\hline
\end{tabular}
\end{table}

\begin{table}[H]
\centering
\caption{Thống kê cấu trúc server}
\label{tab:stats-server}
\scriptsize
\begin{tabular}{|p{5cm}|p{7cm}|}
\hline
\textbf{Thành phần} & \textbf{Số lượng} \\
\hline
Mongoose Schema (Model) & 17 collections \\
\hline
Controller class & 15 controllers \\
\hline
Route file & 6 files (hơn 60 endpoint) \\
\hline
API endpoint & $\sim$65 endpoint \\
\hline
\end{tabular}
\end{table}

\subsection{Minh họa các chức năng chính}
\label{subsection:4.3.3}

Hình \ref{fig:screen-home} minh họa trang chủ storefront với slider banner,
lưới sản phẩm phân trang và chatbot hỗ trợ ở góc phải dưới màn hình.

\begin{figure}[H]
    \centering
    \includegraphics[width=0.9\textwidth]{Hinhve/screen_home.png}
    \caption{Giao diện trang chủ ALIBARBIE}
    \label{fig:screen-home}
\end{figure}

Hình \ref{fig:screen-product} minh họa trang chi tiết sản phẩm với bộ chọn
màu sắc và kích cỡ theo biến thể, hiển thị tồn kho và điều chỉnh giá theo
size được chọn.

\begin{figure}[H]
    \centering
    \includegraphics[width=0.9\textwidth]{Hinhve/screen_product.png}
    \caption{Giao diện chi tiết sản phẩm với bộ chọn biến thể}
    \label{fig:screen-product}
\end{figure}

Hình \ref{fig:screen-checkout} minh họa màn hình thanh toán với form địa chỉ
phân cấp (tỉnh/huyện/xã), hiển thị phí vận chuyển tính tự động và lựa chọn
phương thức thanh toán COD hoặc VNPay.

\begin{figure}[H]
    \centering
    \includegraphics[width=0.9\textwidth]{Hinhve/screen_checkout.png}
    \caption{Giao diện thanh toán với tính phí vận chuyển động}
    \label{fig:screen-checkout}
\end{figure}

Hình \ref{fig:screen-admin-dashboard} minh họa trang Dashboard quản trị với
biểu đồ doanh thu (Chart.js), thống kê tổng đơn hàng, doanh thu tháng và
số lượng khách hàng mới.

\begin{figure}[H]
    \centering
    \includegraphics[width=0.9\textwidth]{Hinhve/screen_admin_dashboard.png}
    \caption{Giao diện Dashboard quản trị với biểu đồ doanh thu}
    \label{fig:screen-admin-dashboard}
\end{figure}

Hình \ref{fig:screen-admin-return} minh họa màn hình quản lý yêu cầu hoàn
hàng với bộ lọc theo trạng thái và lịch sử thay đổi trạng thái từng yêu cầu.

\begin{figure}[H]
    \centering
    \includegraphics[width=0.9\textwidth]{Hinhve/screen_admin_return.png}
    \caption{Giao diện quản lý yêu cầu hoàn hàng}
    \label{fig:screen-admin-return}
\end{figure}

Hình \ref{fig:screen-admin-role} minh họa màn hình phân quyền vai trò RBAC,
cho phép quản trị viên cấu hình danh sách menu và hành động được phép cho
từng vai trò.

\begin{figure}[H]
    \centering
    \includegraphics[width=0.9\textwidth]{Hinhve/screen_admin_role.png}
    \caption{Giao diện quản lý phân quyền vai trò (RBAC)}
    \label{fig:screen-admin-role}
\end{figure}

%%%%%%%%%%%%%%%%%%%%%%%%%%%%%%%%%%%%%%%%%%
\section{Kiểm thử}
\label{section:4.4}

Hệ thống được kiểm thử theo phương pháp \textbf{kiểm thử hộp đen} (black-box
testing) sử dụng kỹ thuật \textbf{phân hoạch tương đương} (equivalence
partitioning) và \textbf{phân tích giá trị biên} (boundary value analysis).
Kiểm thử API thực hiện qua công cụ Postman; kiểm thử giao diện thực hiện thủ
công trên trình duyệt Chrome phiên bản mới nhất.

\subsection*{Use case 1: Đăng nhập hệ thống}

\begin{table}[H]
\centering
\caption{Các trường hợp kiểm thử chức năng đăng nhập}
\label{tab:test-login}
\scriptsize
\begin{tabular}{|p{0.5cm}|p{3.5cm}|p{3cm}|p{3.5cm}|p{2cm}|}
\hline
\textbf{TC} & \textbf{Dữ liệu đầu vào} & \textbf{Điều kiện} & \textbf{Kết quả mong đợi} & \textbf{Kết quả} \\
\hline
1 & Email hợp lệ, mật khẩu đúng & Tài khoản tồn tại & Đăng nhập thành công, set cookie JWT, chuyển trang & Đạt \\
\hline
2 & Email hợp lệ, mật khẩu sai & Tài khoản tồn tại & HTTP 401, thông báo ``Email Hoặc Mật Khẩu Không Chính Xác'' & Đạt \\
\hline
3 & Email không tồn tại & Bất kỳ mật khẩu & HTTP 401, thông báo lỗi xác thực & Đạt \\
\hline
4 & Email để trống & Không nhập & Form báo lỗi validation trước khi gửi request & Đạt \\
\hline
5 & Mật khẩu để trống & Không nhập & Form báo lỗi validation trước khi gửi request & Đạt \\
\hline
6 & Đăng nhập với ``Ghi nhớ đăng nhập'' & Tài khoản hợp lệ & Cookie có maxAge 30 ngày thay vì 1 giờ & Đạt \\
\hline
7 & Email hợp lệ, mật khẩu đúng & Tài khoản có role=admin & Chuyển hướng về trang \texttt{/admin} & Đạt \\
\hline
\end{tabular}
\end{table}

\subsection*{Use case 2: Thêm sản phẩm vào giỏ hàng và đặt hàng}

\begin{table}[H]
\centering
\caption{Các trường hợp kiểm thử chức năng giỏ hàng và đặt hàng}
\label{tab:test-cart}
\scriptsize
\begin{tabular}{|p{0.5cm}|p{3.5cm}|p{3cm}|p{3.5cm}|p{2cm}|}
\hline
\textbf{TC} & \textbf{Dữ liệu đầu vào} & \textbf{Điều kiện} & \textbf{Kết quả mong đợi} & \textbf{Kết quả} \\
\hline
8 & Chọn sản phẩm, size, màu và nhấn ``Thêm vào giỏ'' & Chưa có sản phẩm trong giỏ & Sản phẩm được thêm vào \break Redux store và localStorage & Đạt \\
\hline
9 & Thêm cùng sản phẩm, cùng size/màu lần 2 & Đã có trong giỏ & Không thêm trùng, giữ nguyên giỏ hàng & Đạt \\
\hline
10 & Nhấn thanh toán khi chưa chọn size & Sản phẩm có biến thể & Thông báo yêu cầu chọn size trước & Đạt \\
\hline
11 & Nhập địa chỉ đầy đủ, chọn COD & Đơn hàng hợp lệ & Đơn hàng được tạo, email xác nhận gửi đến khách, giỏ hàng xóa & Đạt \\
\hline
12 & Nhập địa chỉ đầy đủ, chọn VNPay & Đơn hàng hợp lệ & Chuyển hướng đến cổng VNPay & Đạt \\
\hline
13 & Callback VNPay thành công (vnp\_ResponseCode=00) & Giao dịch hợp lệ & Đơn hàng cập nhật \break statusPayment=true, chuyển về trang \texttt{/thanks} & Đạt \\
\hline
14 & Callback VNPay thất bại (vnp\_ResponseCode $\neq$ 00) & Giao dịch thất bại & Đơn hàng giữ \break statusPayment=false, thông báo thất bại & Đạt \\
\hline
15 & Nhập mã coupon hợp lệ & Coupon còn hạn & Hiển thị số tiền được giảm trong tóm tắt giỏ hàng & Đạt \\
\hline
16 & Nhập mã coupon hết hạn & Coupon quá hạn & Thông báo ``Mã giảm giá không hợp lệ'' & Đạt \\
\hline
\end{tabular}
\end{table}

\subsection*{Use case 3: Yêu cầu hoàn hàng}

\begin{table}[H]
\centering
\caption{Các trường hợp kiểm thử chức năng hoàn hàng}
\label{tab:test-return}
\scriptsize
\begin{tabular}{|p{0.5cm}|p{3.5cm}|p{3cm}|p{3.5cm}|p{2cm}|}
\hline
\textbf{TC} & \textbf{Dữ liệu đầu vào} & \textbf{Điều kiện} & \textbf{Kết quả mong đợi} & \textbf{Kết quả} \\
\hline
17 & Gửi yêu cầu với order\_id hợp lệ, lý do hợp lệ & Đơn hàng tồn tại & Tạo return\_request với status PENDING\_REVIEW & Đạt \\
\hline
18 & Nhân viên cập nhật status lên CONTACTING & Role staff & Trạng thái cập nhật thành công, audit log ghi nhận & Đạt \\
\hline
19 & Nhân viên staff cố phê duyệt yêu cầu & Role staff (không đủ quyền) & HTTP 403, thông báo ``Không có quyền'' & Đạt \\
\hline
20 & Quản lý duyệt yêu cầu (approve) & Role manager & Status = APPROVED, email thông báo gửi khách & Đạt \\
\hline
21 & Admin xử lý hoàn tiền & Role admin,\break refund\_amount >0 & surplus của user tăng đúng giá trị, status = REFUNDED & Đạt \\
\hline
\end{tabular}
\end{table}

\subsection*{Use case 4: Đăng ký tài khoản}

\begin{table}[H]
\centering
\caption{Các trường hợp kiểm thử chức năng đăng ký}
\label{tab:test-register}
\scriptsize
\begin{tabular}{|p{0.5cm}|p{3.5cm}|p{3cm}|p{3.5cm}|p{2cm}|}
\hline
\textbf{TC} & \textbf{Dữ liệu đầu vào} & \textbf{Điều kiện} & \textbf{Kết quả mong đợi} & \textbf{Kết quả} \\
\hline
22 & Họ tên, email, mật khẩu, xác nhận mật khẩu hợp lệ & Email chưa tồn tại & Tạo \texttt{ModelUser} mới, mật khẩu được băm bcrypt, chuyển sang trang đăng nhập & Đạt \\
\hline
23 & Đăng ký với email đã tồn tại & Email đã có trong hệ thống & HTTP 403, thông báo ``Người Dùng Đã Tồn Tại'' & Đạt \\
\hline
24 & Email chứa chữ hoa & Bất kỳ & Form báo lỗi validation, không gửi request & Đạt \\
\hline
25 & Mật khẩu và xác nhận mật khẩu không khớp & Bất kỳ & Form báo lỗi validation, không gửi request & Đạt \\
\hline
26 & Để trống họ tên hoặc email hoặc mật khẩu & Không nhập & Form báo lỗi validation tương ứng từng trường & Đạt \\
\hline
\end{tabular}
\end{table}

\subsection*{Use case 5: Quản lý sản phẩm và biến thể}

\begin{table}[H]
\centering
\caption{Các trường hợp kiểm thử chức năng quản lý sản phẩm}
\label{tab:test-product}
\scriptsize
\begin{tabular}{|p{0.5cm}|p{3.5cm}|p{3cm}|p{3.5cm}|p{2cm}|}
\hline
\textbf{TC} & \textbf{Dữ liệu đầu vào} & \textbf{Điều kiện} & \textbf{Kết quả mong đợi} & \textbf{Kết quả} \\
\hline
27 & Tên, giá, danh mục, ảnh sản phẩm hợp lệ & Role admin/manager & \texttt{AddProduct} lưu \break \texttt{ModelProducts} thành công, ghi audit log & Đạt \\
\hline
28 & Thêm sản phẩm thiếu tên hoặc giá & Role admin/manager & Server trả lỗi validation, không lưu bản ghi & Đạt \\
\hline
29 & Sửa giá và mô tả sản phẩm có sẵn & Sản phẩm tồn tại & \texttt{EditProduct} cập nhật thành công, audit log ghi \texttt{data\_before} và \break \texttt{data\_after} & Đạt \\
\hline
30 & Xóa một sản phẩm & Sản phẩm tồn tại & \texttt{DeleteProduct} xóa thành công, audit log ghi nhận hành động & Đạt \\
\hline
31 & Thêm biến thể màu sắc mới cho sản phẩm & Sản phẩm tồn tại & \texttt{AddVariant} lưu thành công, hiển thị thêm lựa chọn màu & Đạt \\
\hline
32 & Thêm size với số lượng tồn kho cho biến thể & Biến thể tồn tại & \texttt{AddVariantSize} lưu thành công, tồn kho hiển thị đúng ở trang chi tiết & Đạt \\
\hline
33 & Nhân viên (role staff) gọi trực tiếp \texttt{POST /api/addproduct} & Role staff (không đủ quyền) & HTTP 403, không tạo sản phẩm & Đạt \\
\hline
\end{tabular}
\end{table}

\subsection*{Use case 6: Quản lý danh mục sản phẩm}

\begin{table}[H]
\centering
\caption{Các trường hợp kiểm thử chức năng quản lý danh mục}
\label{tab:test-category}
\scriptsize
\begin{tabular}{|p{0.5cm}|p{3.5cm}|p{3cm}|p{3.5cm}|p{2cm}|}
\hline
\textbf{TC} & \textbf{Dữ liệu đầu vào} & \textbf{Điều kiện} & \textbf{Kết quả mong đợi} & \textbf{Kết quả} \\
\hline
34 & Tên danh mục hợp lệ & Role admin/manager & \texttt{AddCategory} lưu thành công, hiển thị trong bộ lọc danh mục & Đạt \\
\hline
35 & Để trống tên danh mục & Role admin/manager & Server trả lỗi validation, không lưu bản ghi & Đạt \\
\hline
36 & Sửa tên danh mục đã tồn tại & Danh mục tồn tại & \texttt{EditCategory} cập nhật thành công, sản phẩm thuộc danh mục vẫn hiển thị đúng & Đạt \\
\hline
37 & Xóa một danh mục & Role admin \break (chỉ \texttt{ADMIN\_ONLY}) & \texttt{DeleteCategory} xóa thành công, danh mục không còn trong bộ lọc & Đạt \\
\hline
\end{tabular}
\end{table}

\subsection*{Use case 7: Quản lý vòng đời đơn hàng}

\begin{table}[H]
\centering
\caption{Các trường hợp kiểm thử chức năng quản lý đơn hàng}
\label{tab:test-order}
\scriptsize
\begin{tabular}{|p{0.5cm}|p{3.5cm}|p{3cm}|p{3.5cm}|p{2cm}|}
\hline
\textbf{TC} & \textbf{Dữ liệu đầu vào} & \textbf{Điều kiện} & \textbf{Kết quả mong đợi} & \textbf{Kết quả} \\
\hline
38 & Nhân viên xem danh sách toàn bộ đơn hàng & Role staff/manager/admin & \texttt{GetDataOrder} trả về danh sách kèm sản phẩm từ \texttt{ModelOrderItem} & Đạt \\
\hline
39 & Cập nhật trạng thái đơn hàng từ ``chờ xác nhận'' sang ``đang giao'' & Đơn hàng tồn tại & \texttt{EditOrder} cập nhật \break \texttt{statusOrder} thành công & Đạt \\
\hline
40 & Cập nhật trạng thái đơn hàng thành ``đã giao'' & Đơn hàng tồn tại & \texttt{statusOrder = true}, gửi email xác nhận, đủ điều kiện để khách yêu cầu hoàn hàng & Đạt \\
\hline
41 & Cập nhật trạng thái với \break \texttt{order\_id} không tồn tại & ID không hợp lệ & HTTP 400/404, không thay đổi dữ liệu & Đạt \\
\hline
\end{tabular}
\end{table}

\subsection*{Use case 8: Phân quyền nhân viên (RBAC)}

\begin{table}[H]
\centering
\caption{Các trường hợp kiểm thử chức năng phân quyền}
\label{tab:test-rbac}
\scriptsize
\begin{tabular}{|p{0.5cm}|p{3.5cm}|p{3cm}|p{3.5cm}|p{2cm}|}
\hline
\textbf{TC} & \textbf{Dữ liệu đầu vào} & \textbf{Điều kiện} & \textbf{Kết quả mong đợi} & \textbf{Kết quả} \\
\hline
42 & Tạo tài khoản nhân viên mới với vai trò \texttt{staff} & Role admin/manager & \texttt{CreateUser} tạo thành công, audit log ghi nhận & Đạt \\
\hline
43 & Đổi vai trò nhân viên từ \texttt{staff} sang \texttt{manager} & Role admin & \texttt{UpdateUserRole} cập nhật ngay lập tức, có hiệu lực với request tiếp theo & Đạt \\
\hline
44 & Tạo vai trò mới với danh sách \texttt{menus}/\texttt{actions} tùy chỉnh & Role admin & \texttt{AddRole} lưu thành công vào collection \texttt{role} & Đạt \\
\hline
45 & Role \texttt{manager} gọi \texttt{POST /api/deleterole} & Route chỉ cho \break \texttt{ADMIN\_ONLY} & HTTP 403, không xóa vai trò & Đạt \\
\hline
46 & Nhân viên gọi \texttt{GET /api/my-permissions} & Đã đăng nhập & Trả về đúng danh sách quyền và menu theo vai trò hiện tại & Đạt \\
\hline
\end{tabular}
\end{table}

\subsection*{Use case 9: Nhật ký kiểm toán}

\begin{table}[H]
\centering
\caption{Các trường hợp kiểm thử chức năng nhật ký kiểm toán}
\label{tab:test-audit}
\scriptsize
\begin{tabular}{|p{0.5cm}|p{4.5cm}|p{3cm}|p{3.5cm}|p{1cm}|}
\hline
\textbf{TC} & \textbf{Dữ liệu đầu vào} & \textbf{Điều kiện} & \textbf{Kết quả mong đợi} & \textbf{Kết quả} \\
\hline
47 & Xem danh sách nhật ký có phân trang & Role admin/manager & \texttt{GetAuditLogs} trả về đúng \texttt{data} và \texttt{total} & Đạt \\
\hline
48 & Lọc theo \texttt{actor\_email} cụ thể & Role admin/manager & Chỉ trả về các bản ghi của nhân viên được chọn & Đạt \\
\hline
49 & Lọc theo khoảng thời gian \texttt{from}/\texttt{to} & Role admin/manager & Chỉ trả về bản ghi nằm trong khoảng thời gian đã chọn & Đạt \\
\hline
50 & Xem chi tiết một bản ghi theo ID & Bản ghi tồn tại & \texttt{GetAuditLogById} trả về đầy đủ \texttt{data\_before}, \texttt{data\_after}, \texttt{diff} & Đạt \\
\hline
51 & Role \texttt{staff} gọi \texttt{GET /api/audit-logs} & Route chỉ cho  \break \texttt{ADMIN\_MANAGER} & HTTP 403, không truy cập được nhật ký & Đạt \\
\hline
\end{tabular}
\end{table}

\subsection*{Use case 10: Cấu hình phí vận chuyển}

\begin{table}[H]
\centering
\caption{Các trường hợp kiểm thử chức năng cấu hình phí vận chuyển}
\label{tab:test-shipping}
\scriptsize
\begin{tabular}{|p{0.5cm}|p{4.5cm}|p{3cm}|p{3.5cm}|p{1cm}|}
\hline
\textbf{TC} & \textbf{Dữ liệu đầu vào} & \textbf{Điều kiện} & \textbf{Kết quả mong đợi} & \textbf{Kết quả} \\
\hline
52 & Cấu hình \texttt{domestic\_fee}, \texttt{inter\_province\_fee}, \texttt{free\_threshold} & Role admin/manager & \texttt{UpdateConfig} lưu thành công, có hiệu lực ngay & Đạt \\
\hline
53 & Tính phí vận chuyển khi \texttt{sumPrice} $\geq$ \texttt{free\_threshold} & Đã cấu hình ngưỡng miễn phí & \texttt{CalculateFee} trả về phí bằng 0 & Đạt \\
\hline
54 & Tính phí vận chuyển cùng tỉnh/thành với kho hàng & Địa chỉ cùng tỉnh & Trả về đúng \break \texttt{domestic\_fee} đã cấu hình & Đạt \\
\hline
55 & Tính phí vận chuyển khác tỉnh/thành & Địa chỉ khác tỉnh & Trả về đúng \break \texttt{inter\_province\_fee} đã cấu hình & Đạt \\
\hline
56 & Role \texttt{staff} gọi \texttt{POST /api/shipping-config} & Route chỉ cho \break \texttt{ADMIN\_MANAGER} & HTTP 403, không thay đổi cấu hình & Đạt \\
\hline
\end{tabular}
\end{table}

\subsection*{Use case 11: Quản lý mã giảm giá}

\begin{table}[H]
\centering
\caption{Các trường hợp kiểm thử chức năng quản lý mã giảm giá}
\label{tab:test-coupon}
\scriptsize
\begin{tabular}{|p{0.5cm}|p{4.5cm}|p{3cm}|p{3.5cm}|p{1cm}|}
\hline
\textbf{TC} & \textbf{Dữ liệu đầu vào} & \textbf{Điều kiện} & \textbf{Kết quả mong đợi} & \textbf{Kết quả} \\
\hline
57 & Tạo mã giảm giá với \break \texttt{discount\_percent} và \break \texttt{expiry\_date} hợp lệ & Role admin/manager & \texttt{AddCoupon} lưu thành công & Đạt \\
\hline
58 & Sửa \texttt{usage\_limit} của mã giảm giá có sẵn & Mã tồn tại & \texttt{EditCoupon} cập nhật thành công & Đạt \\
\hline
59 & Xóa một mã giảm giá & Role admin \break (chỉ \texttt{ADMIN\_ONLY}) & \texttt{DeleteCoupon} xóa thành công, không còn hiển thị cho khách hàng & Đạt \\
\hline
60 & Khách hàng gọi \texttt{GET /api/available-coupons} & Có mã còn hạn và còn lượt dùng & Chỉ trả về các mã thỏa \break \texttt{usage\_limit > 0} và chưa hết hạn & Đạt \\
\hline
61 & Role \texttt{staff} gọi \texttt{POST /api/deletecoupon} & Route chỉ cho \break \texttt{ADMIN\_ONLY} & HTTP 403, không xóa mã giảm giá & Đạt \\
\hline
\end{tabular}
\end{table}

\subsection*{Use case 12: Danh sách yêu thích}

\begin{table}[H]
\centering
\caption{Các trường hợp kiểm thử chức năng danh sách yêu thích}
\label{tab:test-wishlist}
\scriptsize
\begin{tabular}{|p{0.5cm}|p{3.5cm}|p{3cm}|p{3.5cm}|p{2cm}|}
\hline
\textbf{TC} & \textbf{Dữ liệu đầu vào} & \textbf{Điều kiện} & \textbf{Kết quả mong đợi} & \textbf{Kết quả} \\
\hline
62 & Nhấn yêu thích sản phẩm chưa có trong danh sách & Đã đăng nhập & \texttt{Toggle} thêm sản phẩm vào wishlist, biểu tượng cập nhật trạng thái đã thích & Đạt \\
\hline
63 & Nhấn yêu thích lần thứ hai trên cùng sản phẩm & Sản phẩm đã có trong wishlist & \texttt{Toggle} gỡ sản phẩm khỏi wishlist & Đạt \\
\hline
64 & Kiểm tra trạng thái yêu thích của sản phẩm & Đã đăng nhập & \texttt{CheckLiked} trả về đúng trạng thái true/false & Đạt \\
\hline
65 & Khách chưa đăng nhập nhấn yêu thích & Chưa đăng nhập & Hệ thống yêu cầu đăng nhập trước khi thêm vào wishlist & Đạt \\
\hline
\end{tabular}
\end{table}

\subsection*{Use case 13: Quản lý blog}

\begin{table}[H]
\centering
\caption{Các trường hợp kiểm thử chức năng quản lý blog}
\label{tab:test-blog}
\scriptsize
\begin{tabular}{|p{0.5cm}|p{3.5cm}|p{3cm}|p{3.5cm}|p{2cm}|}
\hline
\textbf{TC} & \textbf{Dữ liệu đầu vào} & \textbf{Điều kiện} & \textbf{Kết quả mong đợi} & \textbf{Kết quả} \\
\hline
66 & Thêm bài blog với tiêu đề và nội dung hợp lệ & Role admin/manager & \texttt{AddBlog} lưu thành công, hiển thị ở trang chủ & Đạt \\
\hline
67 & Xóa một bài blog & Bài blog tồn tại & \texttt{DeleteBlog} xóa thành công, không còn hiển thị & Đạt \\
\hline
68 & Khách hàng xem danh sách blog & Không yêu cầu đăng nhập & \texttt{GetBlog} trả về đúng danh sách bài đã đăng & Đạt \\
\hline
\end{tabular}
\end{table}

\subsection*{Use case 14: Chatbot hỗ trợ thời gian thực}

\begin{table}[H]
\centering
\caption{Các trường hợp kiểm thử chức năng chatbot}
\label{tab:test-chatbot}
\scriptsize
\begin{tabular}{|p{0.5cm}|p{3.5cm}|p{3cm}|p{3.5cm}|p{2cm}|}
\hline
\textbf{TC} & \textbf{Dữ liệu đầu vào} & \textbf{Điều kiện} & \textbf{Kết quả mong đợi} & \textbf{Kết quả} \\
\hline
69 & Khách hàng gửi tin nhắn qua sự kiện \texttt{sendMessage} & Kết nối Socket.io đã thiết lập & Nhân viên hỗ trợ nhận tin nhắn theo thời gian thực, độ trễ dưới 100ms & Đạt \\
\hline
70 & Mất kết nối mạng giữa chừng & Đang trong phiên chat & Client Socket.io tự kết nối lại khi mạng phục hồi & Đạt \\
\hline
71 & Nhân viên trả lời tin nhắn của khách & Phiên chat đang mở & Khách hàng nhận được phản hồi tức thời trên giao diện & Đạt \\
\hline
\end{tabular}
\end{table}

\subsubsection*{Tổng kết kiểm thử}

Tổng cộng 71 trường hợp kiểm thử được thực hiện trên 14 nhóm chức năng chính
của hệ thống, bao gồm các kịch bản hợp lệ (positive) và không hợp lệ
(negative), trong đó có các trường hợp kiểm tra phân quyền truy cập API theo
vai trò (TC-19, TC-33, TC-45, TC-51, TC-56, TC-61). Kết quả: \textbf{71/71
trường hợp đạt} (100\%).

%%%%%%%%%%%%%%%%%%%%%%%%%%%%%%%%%%%%%%%%%%
\section{Triển khai}
\label{section:4.5}

\subsubsection*{Mô hình triển khai thử nghiệm}

Hệ thống được triển khai thử nghiệm trên môi trường phát triển cục bộ
(localhost) với cấu hình như mô tả trong Bảng \ref{tab:deploy}.

\begin{table}[H]
\centering
\caption{Cấu hình môi trường triển khai thử nghiệm}
\label{tab:deploy}
\scriptsize
\begin{tabular}{|p{4cm}|p{9cm}|}
\hline
\textbf{Thành phần} & \textbf{Cấu hình} \\
\hline
Phần cứng & CPU Intel Core i5, RAM 8 GB, SSD 256 GB \\
\hline
Hệ điều hành & Windows 11 Pro \\
\hline
Node.js Server & Cổng 5000, khởi động bằng Nodemon \\
\hline
React Dev Server & Cổng 3000, khởi động bằng \texttt{react-scripts start} \\
\hline
MongoDB & Cổng 27017, cài đặt cục bộ (MongoDB Community 7.0) \\
\hline
VNPay & Môi trường Sandbox (tài khoản thử nghiệm) \\
\hline
Email (SMTP) & Google SMTP qua OAuth2 (tài khoản Gmail thử nghiệm) \\
\hline
\end{tabular}
\end{table}

\subsubsection*{Quy trình khởi động hệ thống}

Hệ thống được khởi động theo hai bước:

\begin{enumerate}
  \item Khởi động máy chủ API:
        \texttt{cd server \&\& npm start}
        (Nodemon theo dõi thay đổi và tự khởi động lại).
  \item Khởi động ứng dụng React:
        \texttt{cd client \&\& npm start}
        (Trình duyệt tự mở tại \texttt{http://localhost:3000}).
\end{enumerate}

Biến môi trường được cấu hình qua file \texttt{.env} phía server, bao gồm:
chuỗi kết nối MongoDB (\texttt{MONGO\_URI}), khóa bí mật JWT
(\texttt{JWT\_SECRET}), thông tin VNPay sandbox (\texttt{VNP\_TMN\_CODE},
\texttt{VNP\_HASH\_SECRET}) và thông tin OAuth2 Google (\texttt{GOOGLE\_CLIENT\_ID},
\texttt{GOOGLE\_CLIENT\_SECRET}).

\subsubsection*{Kết quả triển khai thử nghiệm}

Hình \ref{fig:deploy-arch} mô tả kiến trúc triển khai thực tế của hệ thống.

\begin{figure}[H]
    \centering
    % SV bổ sung deployment diagram tại đây
    \includegraphics[width=0.8\textwidth]{Hinhve/deploy_arch.png}
    \caption{Mô hình triển khai hệ thống ALIBARBIE}
    \label{fig:deploy-arch}
\end{figure}

Kết quả kiểm thử triển khai cho thấy:

\begin{itemize}
  \item Thời gian phản hồi API trung bình dưới 200ms cho các truy vấn đọc
        thông thường (lấy danh sách sản phẩm, giỏ hàng), đáp ứng yêu cầu dưới
        2 giây ở mục \ref{section:2.4}.
  \item Giao tiếp Socket.io duy trì kết nối ổn định, độ trễ thông điệp chat
        dưới 100ms trên mạng cục bộ.
  \item Luồng thanh toán VNPay sandbox hoạt động chính xác với tất cả các mã
        phản hồi thử nghiệm.
  \item Giao diện hiển thị đúng trên Chrome, Firefox và Edge; tương thích
        thiết bị di động tại breakpoint 375px (iPhone SE).
\end{itemize}

Chương này đã trình bày thiết kế kiến trúc ba tầng--MVC, thiết kế chi tiết
giao diện, lớp và cơ sở dữ liệu, thống kê kết quả xây dựng, 71 trường hợp
kiểm thử trên 14 nhóm chức năng và kết quả triển khai thử nghiệm của hệ thống
ALIBARBIE. Chương tiếp theo sẽ tổng kết những kết quả đạt được và hướng phát
triển tiếp theo.

%%%%%%%%%%%%%%%%%%%%
\end{document}
